\pagebreak
\section{Метод решения}
Поисковая система --- булев поиск из ЛР7 (индекс) и ЛР8 (поиск) по корпусу из 30\,000 английских статей о туризме (Wikivoyage и Wikipedia) --- первому корпусу ЛР1. Остальные работы сделаны на корпусе в 1\,000\,000 документов, но оценки релевантности в этой работе привязаны к номерам документов первого корпуса, поэтому она оставлена на нём. Выполнен простейший вариант задания: поиск без учёта словоформ. Стемминг применяется и при индексации (в словарь попадают основы слов), и при выполнении запроса (слова запроса приводятся к тем же основам). Бонус за точное совпадение не делался: булев поиск не ранжирует выдачу.

\subsection{Стеммер Портера}
Корпус английский, поэтому взят стеммер Портера (M.\,F.~Porter, \enquote{An algorithm for suffix stripping}, 1980). Он написан самостоятельно по описанию алгоритма, без библиотек: \texttt{lab4/src/porter.h}, около 200 строк C++.

Слово рассматривается как последовательность согласных (C) и гласных (V) букв ($y$ после согласной считается гласной) вида $[C](VC)^m[V]$; $m$ --- мера основы. Суффиксы снимаются за пять шагов, на каждом шаге применяется правило с самым длинным подходящим суффиксом, если основа удовлетворяет условию:
\begin{enumerate}
    \item множественное число и формы на \emph{-ed}, \emph{-ing}: \emph{caresses} $\to$ \emph{caress}, \emph{ponies} $\to$ \emph{poni}, \emph{hopping} $\to$ \emph{hop}, \emph{hiking} $\to$ \emph{hike}; конечная $y$ после гласной в основе $\to$ $i$ (\emph{happy} $\to$ \emph{happi});
    \item двойные суффиксы при $m > 0$: \emph{-ational} $\to$ \emph{-ate}, \emph{-ization} $\to$ \emph{-ize}, \emph{-fulness} $\to$ \emph{-ful} и ещё 17 правил;
    \item \emph{-icate}, \emph{-ative}, \emph{-alize}, \emph{-ical}, \emph{-ful}, \emph{-ness} при $m > 0$;
    \item суффиксы \emph{-al}, \emph{-ance}, \emph{-er}, \emph{-ic}, \emph{-able}, \emph{-ment}, \emph{-ion}, \emph{-ism}, \emph{-ive}, \emph{-ize} и др. снимаются при $m > 1$;
    \item конечная \emph{e} и двойная \emph{ll}.
\end{enumerate}
Пример: \emph{generalization} $\to$ \emph{generalize} (шаг 2) $\to$ \emph{general} (шаг 3) $\to$ \emph{gener} (шаг 4).

Стеммер применяется только к токенам из латинских букв \texttt{a--z} длиннее двух символов; токены с цифрами, диакритикой и не латиницей (\texttt{10am}, \texttt{café}, \texttt{москва}) остаются как есть.

\subsection{Встраивание в поисковую систему}
\begin{itemize}
    \item Индексатор ЛР7 получил ключ \texttt{-s}: каждый токен перед добавлением в словарь заменяется основой. В поле \texttt{flags} заголовка индекса ставится бит~0 \enquote{термы --- основы}. Это расширение формата, для которого поле \texttt{flags} и было зарезервировано в ЛР7: старые индексы (флаг 0) читаются как раньше.
    \item Поиск ЛР8 при загрузке индекса читает флаг и, если он стоит, приводит к основе каждое слово запроса после токенизации. Операторы и разбор запроса не меняются. В строке \texttt{parsed:} видно, во что превратился запрос: \texttt{skiing resorts} $\to$ \texttt{ski \&\& resort}.
    \item Утилита \texttt{lab4/build/stem} печатает основу для каждого токена текста --- для тестов и анализа.
\end{itemize}
Индекс со стеммингом строится отдельно (\texttt{lab4/data/index-stem.bin}); индекс без стемминга из ЛР7 остаётся, поэтому обе системы можно сравнить на одних и тех же запросах.

\subsection{Оценка качества}
\begin{itemize}
    \item \textbf{Запросы.} 14 запросов туристической тематики (\texttt{data/queries.tsv}), у каждого записана информационная потребность --- критерий релевантности. Запросы 1--10 содержат словоформы, на которых стемминг должен помочь (множественное число, \emph{-ing}): \texttt{cheap hostels}, \texttt{skiing resorts}, \texttt{snorkeling coral reefs thailand} и др. Запросы 11--14 подобраны там, где стеммер Портера склеивает слова с разным смыслом: \texttt{travel news} (\emph{news} $\to$ \emph{new}), \texttt{general hospital} (\emph{gener}, \emph{hospit}), \texttt{organ concerts} (\emph{organ} = \emph{organization}), \texttt{university museums} (\emph{univers} = \emph{universe}).
    \item \textbf{Метрика.} Булев поиск выдаёт документы в порядке номеров (сначала Wikivoyage), пользователь видит первую страницу. Поэтому мера --- точность первых 10 результатов, P@10: доля релевантных среди первых 10, если найдено меньше 10, недостающие считаются нерелевантными. Для полноты --- число найденных документов и грубая оценка числа найденных релевантных: найдено $\times$ P@10.
    \item \textbf{Разметка.} Пул --- объединение первых 10 результатов обеих систем по каждому запросу, 225 документов (\texttt{results/pool.txt}: заголовок и фрагменты текста со словами запроса). Оценки 0/1 ставились вслепую: в файле для разметки не было видно, какая система нашла документ, при нехватке фрагментов читался полный текст. Оценки с обоснованием каждой --- \texttt{data/qrels.tsv}.
\end{itemize}
